\BeleskeID{09_2-s048}
% element: 09_2-s048:d01
% element: 09_2-s048:e05
% element: 09_2-s048:e06
Radi skaliranja i prelaska sa jedne tehnologije na drugu uvedena je normalizovana dimenzija $\lambda$. Dimenzije se zadaju celobrojnim umnošcima tog parametra, pri čemu $2\lambda$ predstavlja minimum. Zato se neke dimenzije mogu ostvariti i kao polovina minimalne dimenzije.

Pri označavanju tranzistora najčešće se podrazumeva zajednička, često minimalna dužina kanala, pa se uz simbol ispisuju samo širine. Jedinični CMOS invertor bira jednake dužine P i N kanala, dok širine zadovoljavaju neki od izvedenih odnosa ili njihov kompromis, najčešće 2:1. Takav invertor zauzima najmanje prostora i ima minimalne kapacitivnosti u posmatranom izboru.

Odnos za veće tranzistore treba zaokružiti tek pošto se izračuna potrebna širina. Za odnos 2.376 i $W_n=10$ račun daje 23.76, pa se bira 24, a ne 20. Prerano zaokruživanje samog odnosa na dva prenelo bi znatno veću grešku na dimenzionisanje.
